\begin{remark}[7.5.8]
\label{IV.2.s7.5.8-fr}
Sería interesante determinar si, en (7.5.7), se puede sustituir
la hipótesis de que $k$ es perfecto o algebraicamente cerrado por la de que
$\operatorname{Spec}(A)$ o $\operatorname{Spec}(B)$ es \emph{geométricamente
íntegro} sobre $k$, o al menos por la hipótesis de que $A$ (por ejemplo) sea
íntegro y contenga un subanillo $A_0$ isomorfo a un anillo de series
formales $k[[T_1,\ldots,T_n]]$, tal que el cuerpo de fracciones $K$ de $A$ sea
\emph{separable} sobre el cuerpo de fracciones $K_0$ de $A_0$ (puede demostrarse que esta
condición implica que $A$ es geométricamente íntegro sobre $k$, y que
ambas propiedades son equivalentes cuando $[k:k^p]$ es finito, siendo $p$
el exponente característico de $k$). También sería deseable desarrollar
variantes de (7.5.6) y (7.5.7) en las que se debilitara la hipótesis de finitud
de los cuerpos residuales de $A$ y $B$, suponiendo, por ejemplo, que sólo uno de
ellos sea una extensión finita de $k$, mientras que el otro sea arbitrario.
\end{remark}
